\documentclass{article}
\usepackage{amsmath}
\usepackage{amssymb}
\usepackage{array}
\usepackage{booktabs}
\usepackage{textcomp}
\usepackage{graphicx}
\usepackage{tikz}
\usepackage{graphics}
\usepackage{pdfpages}

\usepackage{amsthm}
\newtheorem{theorem}{Theorem}[section]

\newtheorem{lemma}[theorem]{Lemma}



\newcommand{\coNP}{\textrm{co-}\mathbb{NP}}

\title{cse 242 hw1}
\author{Ian Chi}
\date{\today}
\begin{document}
\maketitle

\section{}

\begin{enumerate}
    \item 
    \begin{enumerate}
        \item the characteristic polynomial for this matrix \(\Sigma\) is 
        \begin{equation}
            \lambda^2-41\lambda+175
        \end{equation}
        which results in eigenvalues of 
        
        \begin{equation}
            \frac{1}{2}(41\pm 3\sqrt{109})
        \end{equation}

        setting \(\left( \Sigma-\lambda 1 \right) v =0\), solving for \(v\) gives the eigenvectors of \(\lambda\). this results in \((\frac{1}{10}(3\pm\sqrt{109}),1)\)

        we can check if 

        \begin{equation}
            \begin{pmatrix}
                25 & 15\\
                15&16
            \end{pmatrix}
            \begin{pmatrix}
                \frac{1}{10}(3\pm\sqrt{109})\\
                1
            \end{pmatrix}
             = \frac{1}{2}(41+3\sqrt{109})\begin{pmatrix}
                \frac{1}{10}(3\pm\sqrt{109})\\
                1
            \end{pmatrix}
        \end{equation}

        and we find that this equation holds. 

        \begin{equation}
            \begin{split}
            \begin{pmatrix}
                25 & 15\\
                15&16
            \end{pmatrix}
            \begin{pmatrix}
                \frac{1}{10}(3\pm\sqrt{109})\\
                1
            \end{pmatrix} &= \begin{pmatrix}
                15+\frac{5}{2}(3\pm\sqrt{109})\\
                16+\frac{3}{2}(3\pm\sqrt{109})\\
            \end{pmatrix}\\
            &=\frac{1}{2}(41+3\sqrt{109})\begin{pmatrix}
                \frac{1}{10}(3\pm\sqrt{109})\\
                1
            \end{pmatrix}
            \end{split}
        \end{equation}

        \item a plot of the pdf would have 1 std extend 4 units in 1 axis and 5 in the other. the level curve would be an ellipse with some positive correlation between the two axis since the cov is positive. 
        \item the direction that preserves the most variance is the one corresponding to the largest eigenvector of the matrix. 
        
        the fraction of variance preserved would be \(\lambda/41\approx 0.88\) where 41 is the trace of the cov matrix. 
    \end{enumerate}
    \item \begin{equation}
        \begin{split}
            Tr(AB)&= \sum_i^m(AB)_{ii}\\
            &= \sum_i^{m}\sum_j^n (A_{ij}B_{ji}) \\
            &= \sum_j^{n}\sum_i^m (A_{ij}B_{ji}) \\
            &= \sum_j^n(BA)_{jj}\\
            &= Tr(BA)\\
        \end{split}
    \end{equation}
    \item 
    \begin{equation}
        \begin{split}
            E[x^TAx]&=E[(\mu+(x-\mu))^TA(\mu+(x-\mu))]\\
            &= E[(\mu^TA\mu) + (\mu^TA(x-\mu)) \\
            &\quad+ ((x-\mu)^TA\mu) + ((x-\mu)^TA((x-\mu)))]\\
            &= (\mu^TA\mu) + E[(\mu^TA(x-\mu)) \\
            &\quad+ ((x-\mu)^TA\mu) + ((x-\mu)^TA((x-\mu)))]\\
            &= (\mu^TA\mu) + E[(\mu^TA(x-\mu)) \\
            &\quad+ ((x-\mu)^TA\mu) + ((x-\mu)^TA((x-\mu)))]\\
            &= (\mu^TA\mu) + E[((x-\mu)^TA((x-\mu)))]\\
            &= (\mu^TA\mu) + E[A((x-\mu)^T((x-\mu)))]\\
            &= (\mu^TA\mu) + tr(A\epsilon)
            % &= (\mu^TA\mu) + E[((x-\mu)^TA\mu) + ((x)^TA((x-\mu)))]\\
        \end{split}
    \end{equation}
    \item 

    \begin{equation}
        E[X]=\int_{-1}^1 \frac{x}{2} dx = 0
    \end{equation}
    \begin{equation}
        E[X^2]=\int_{-1}^1 \frac{x^2}{2} dx = \frac{1}{3}
    \end{equation}
    \begin{equation}
        E[X^3]=\int_{-1}^1 \frac{x^3}{2} dx = 0
    \end{equation}
    \begin{equation}
        E[XY]=\int_{-1}^1 \frac{x^3}{2} dx = 0
    \end{equation}
    Cov is however, 

    \begin{equation}
        E[X^3]-E[X]E[X^2]=0
    \end{equation}

    since cov is zero, so is correlation but XY are obviously dependent. 
\end{enumerate}

\section{}
\begin{enumerate}
    \item \begin{enumerate}
        \item \begin{equation}\begin{split}
                P(d|+)=\frac{P(+|d)P(d)}{P(+)} &= 
                P(d|+)=\frac{P(+|d)P(d)}{P(+ | d) \cdot P(d) + P(+ | \neg d) \cdot P(\neg d)} \\
                &=\frac{1}{100}
        \end{split}
        \end{equation}
        $$P(+ | d) \cdot P(d) = 0.99 \cdot 0.0001 = 0.000099$$
        $$P(+) = (0.99 \cdot 0.0001) + (0.01 \cdot 0.9999)$$
        $$P(D | +) = \frac{0.000099}{0.010098}\approx 0.0098$$
        \item this statement doesnt work because it doesnt take in account of the prior probability of having the disease.
        \item $$P(d | +) = \frac{P(+ | d) \cdot P(d)_{new}}{P(+)_{new}}$$
        $$0.99 \cdot 0.0098039 \approx 0.0097059$$
        $$P(+) = (0.99 \cdot 0.0098039) + (0.01 \cdot 0.9901961) \approx 0.0196$$
        $$P(d | +_{second}) = \frac{0.0097059}{0.0196079} \approx 0.495$$
        \item $$\frac{P(+|d) \cdot P(d)}{P(+|d) \cdot P(d) + (1 - s) \cdot P(\neg d)} \ge 0.5$$
        \(s \ge 0.99990099...\)
    \end{enumerate}
    \item \begin{enumerate}
        \item the prosecutor is miscalculating the prior. since the prior is known, the statement that the chance is exactly .99 of being guilty is wrong
        \item the defender is not acknowledging that having this evidence changes the prior. future evidence will build on this new prior instead of the old one. 
        \item you would need to know the prior probability of the defendant being guilty. 
    \end{enumerate}

    \item \begin{enumerate}
        \item maximize \begin{equation}
            P^3((1-P)^2)
        \end{equation}
        we get the max at .6
        \item instead we maximize \begin{equation}
            P^3((1-P)^2)2P
        \end{equation}
        which gives \(P=\frac{2}{3}\)
        \item \(a=b=1\) means there were two flips, one heads one tails. when they equal two, there were 4 flips w 2 heads 2 tails. as \(a+b\) gets large, it means that there were many flips before and that you are very certain to the probability of heads. \(\hat \theta_{MAP}\) therefore does not change much given an extra sample point or two. 
        
    \end{enumerate}
    \item $$L(\mu, \sigma^2) = \prod_{i=1}^{n} \frac{1}{\sqrt{2\pi\sigma^2}} \exp\left( -\frac{(x_i - \mu)^2}{2\sigma^2} \right)$$$$L(\mu, \sigma^2) = \left( 2\pi\sigma^2 \right)^{-n/2} \exp\left( -\frac{1}{2\sigma^2} \sum_{i=1}^{n} (x_i - \mu)^2 \right)$$
        
        we can take log to make derivatives easier while preserving maximums. 

        $$l(\mu, \sigma^2) = \ln \left[ \left( 2\pi\sigma^2 \right)^{-n/2} \exp\left( -\frac{1}{2\sigma^2} \sum_{i=1}^{n} (x_i - \mu)^2 \right) \right]$$

        solving for \(\mu,\) we can take \(\sigma^2\) to be constant,      
        $$\frac{\partial l}{\partial \mu} = \frac{\partial}{\partial \mu} \left( -\frac{1}{2\sigma^2} \sum_{i=1}^{n} (x_i - \mu)^2 \right)=\frac{1}{\sigma^2} \left( \sum_{i=1}^{n} x_i - n\mu \right) = 0$$

        gives \(\mu = 4\)

        $$\frac{\partial l}{\partial \sigma^2} = -\frac{n}{2} \frac{1}{\sigma^2} + \frac{1}{2(\sigma^2)^2} \sum_{i=1}^{n} (x_i - \mu)^2$$

        implies $$-n\sigma^2 + \sum_{i=1}^{n} (x_i - \mu)^2 = 0 \implies \hat{\sigma}^2 = \frac{1}{n} \sum_{i=1}^{n} (x_i - \mu)^2$$
        $$\hat{\sigma}^2 = \frac{(1-4)^2 + (3-4)^2 + (5-4)^2 + (7-4)^2}{4}=5$$
        \item 
        \begin{enumerate}
        \item Likelihood: $P(x_1 | \mu) = \frac{1}{\sqrt{2\pi \sigma^2}} \exp\left(-\frac{(x_1 - \mu)^2}{2\sigma^2}\right)$

        Prior: $P(\mu) = \frac{1}{\sqrt{2\pi \tau^2}} \exp\left(-\frac{(\mu - \mu_0)^2}{2\tau^2}\right)$ where $\mu_0 = 0$ and $\tau^2 = 1$.


        $$\log P(\mu | x_1) \propto -\frac{(5 - \mu)^2}{2(4)} - \frac{\mu^2}{2(1)}$$

        $$\frac{d}{d\mu} \left[ -\frac{(5 - \mu)^2}{8} - \frac{\mu^2}{2} \right] = \frac{5 - \mu}{4} - \mu = 0$$
        this gives \(\hat \mu = 1\)
        \item we instead get $$\log P(\mu | X) \propto -\sum_{i=1}^n \frac{(x_i - \mu)^2}{2\sigma^2} - \frac{(\mu - \mu_0)^2}{2\tau^2}$$
        $$\frac{d}{dx} \log P(\mu | X) = \frac{1}{\sigma^2} \sum_{i=1}^n (x_i - \mu) - \frac{\mu - \mu_0}{\tau^2} = 0$$
        we get \(\hat \mu = \frac{n}{n+4}\bar x\) where \(\bar x\) represents the mean of the samples. 

        \item the first term in the \(\hat \mu\) in part \(b\) goes to 1 which means it converges to the mean of \(x\) as \(n\to \infty\).
        \end{enumerate}
        
        \item \begin{enumerate}
        \item \begin{equation}
            \begin{split}
                Var(X)&=\frac{\sum ( x-\mu)^2}{N}\\
                &=\frac{\sum x^2 - 2x\mu + \mu^2}{N}\\
                &=\frac{\sum x^2 - \sum 2x\mu + \mu^2}{N}\\
                &=\frac{\sum x^2}{N} - \frac{2\mu \sum x}{N} + \mu^2\\
                &=\frac{\sum x^2}{N} - {2\mu \mu }{} + \mu^2\\
                &=E[X^2]-E[X]^2
            \end{split}
        \end{equation}
        \item 

        \begin{equation}
            \begin{split}
                Var(aX+b) &= E[(aX+b)^2] - E[aX+b]^2\\
                &= E[(a^2X^2 + 2aXb+b^2)] - (aE[X]+b)^2\\
                &= E[(a^2X^2 + 2aXb+b^2)] \\
                &\quad - [(aE[X])^2 + 2aE[X]b +b^2]\\
                &= E[(a^2X^2 )] - (aE[X])^2 \\
                &= a^2 Var(X)\\
            \end{split}
        \end{equation}
        \item 

        \item
        \begin{align}
        Var(X+Y) &= E[(X+Y)^2] - E[X+Y]^2\\
        &= E[X^2+2XY+Y^2] - (E[X]+E[Y])^2\\
        &= E[X^2+2XY+Y^2] - (E[X]+E[Y])^2\\
        &= E[X^2+2XY+Y^2] \\
        &\quad - (E[X]^2 + 2E[X]E[Y] + E[Y]^2)\nonumber\\
        &= E[X^2]+E[2XY]+E[Y^2]\label{ind1} \\
        &\quad - (E[X]^2 + 2E[X]E[Y] + E[Y]^2)\nonumber\\
        &= E[X^2]+2E[X]E[Y]+E[Y^2] \label{ind2}\\
        &\quad - (E[X]^2 + 2E[X]E[Y] + E[Y]^2)\nonumber\\
        &= E[X^2]+E[Y^2] - (E[X]^2  + E[Y]^2)\\ 
        &= Var(X)+Var(Y)
        \end{align}
        note the lines~\ref{ind1} and~\ref{ind2} are the lines that require independence. \(E[XY]=E[X]E[Y]\)


        \end{enumerate}
         
\end{enumerate}


\end{document}